\documentclass[11pt,a4paper]{article}

% ── Packages ──────────────────────────────────────────────────────────────────
\usepackage[utf8]{inputenc}
\usepackage[T1]{fontenc}
\usepackage{lmodern}
\usepackage[margin=2.5cm]{geometry}
\usepackage{hyperref}
\usepackage{xcolor}
\usepackage{listings}
\usepackage{enumitem}
\usepackage{graphicx}
\usepackage{fancyhdr}
\usepackage{titlesec}
\usepackage{tcolorbox}

% ── Colors ────────────────────────────────────────────────────────────────────
\definecolor{cmdBg}{HTML}{F5F5F5}
\definecolor{cmdFrame}{HTML}{CCCCCC}
\definecolor{linkBlue}{HTML}{0969DA}
\definecolor{warnOrange}{HTML}{D29922}
\definecolor{tipGreen}{HTML}{1A7F37}

% ── Hyperlinks ────────────────────────────────────────────────────────────────
\hypersetup{
    colorlinks=true,
    linkcolor=linkBlue,
    urlcolor=linkBlue,
    pdftitle={ro.uvt.ri -- Developer Guide},
    pdfauthor={Team Q-UW9}
}

% ── Listings (terminal commands) ──────────────────────────────────────────────
\lstdefinestyle{terminal}{
    backgroundcolor=\color{cmdBg},
    frame=single,
    rulecolor=\color{cmdFrame},
    basicstyle=\ttfamily\small,
    breaklines=true,
    columns=fullflexible,
    keepspaces=true,
    xleftmargin=4pt,
    xrightmargin=4pt,
    aboveskip=8pt,
    belowskip=8pt
}
\lstset{style=terminal}

% ── Colored boxes ─────────────────────────────────────────────────────────────
\newtcolorbox{warningbox}{
    colback=warnOrange!8,
    colframe=warnOrange!60,
    title={\textbf{Warning}},
    fonttitle=\bfseries\color{warnOrange!80!black},
    sharp corners,
    boxrule=0.5pt
}

\newtcolorbox{tipbox}{
    colback=tipGreen!6,
    colframe=tipGreen!50,
    title={\textbf{Tip}},
    fonttitle=\bfseries\color{tipGreen!80!black},
    sharp corners,
    boxrule=0.5pt
}

% ── Header / Footer ──────────────────────────────────────────────────────────
\pagestyle{fancy}
\fancyhf{}
\fancyhead[L]{\small\texttt{ro.uvt.ri} -- Developer Guide}
\fancyhead[R]{\small Team Q-UW9}
\fancyfoot[C]{\thepage}
\renewcommand{\headrulewidth}{0.4pt}

% ── Title ─────────────────────────────────────────────────────────────────────
\title{%
    \textbf{ro.uvt.ri} \\[4pt]
    \large Developer Onboarding \& Workflow Guide
}
\author{Team Q-UW9}
\date{\today}

% ══════════════════════════════════════════════════════════════════════════════
\begin{document}
\maketitle
\tableofcontents
\newpage

% ──────────────────────────────────────────────────────────────────────────────
\section{Project Overview}
% ──────────────────────────────────────────────────────────────────────────────

This project follows a \textbf{Headless CMS} architecture:

\begin{itemize}[leftmargin=1.5em]
    \item \textbf{Backend (CMS):} WordPress -- used exclusively for content
          management via the REST API. Each developer runs a local instance using Laragon.
    \item \textbf{Frontend:} React (Vite + TypeScript) -- the actual website
          that end-users interact with. It consumes the WordPress REST API.
\end{itemize}

\subsection{Repository Structure}

\begin{lstlisting}
ro.uvt.ri/
|-- .gitignore          # Ignores WP core, node_modules, env files
|-- README.md           # Quick-start instructions
|-- docs/               # This guide and other documentation
|   +-- developer-guide.tex
+-- frontend/           # React application (Vite + TypeScript)
    |-- package.json
    |-- src/
    |   |-- App.tsx
    |   |-- api/wordpress.js
    |   +-- pages/
    +-- vite.config.ts
\end{lstlisting}

\begin{warningbox}
The \texttt{wordpress/} folder on your local machine is \textbf{not tracked by Git}.
Each developer sets up their own local WordPress instance using a local server stack such as Laragon.
\end{warningbox}

% ──────────────────────────────────────────────────────────────────────────────
\section{Prerequisites}
% ──────────────────────────────────────────────────────────────────────────────

Before cloning the repository, make sure the following tools are installed:

\begin{enumerate}[leftmargin=1.5em]
    \item \textbf{Git} --
          \url{https://git-scm.com/downloads}
    \item \textbf{Node.js} (v18 or newer, LTS recommended) --
          \url{https://nodejs.org/}
    \item \textbf{Laragon} (or similar XAMPP/WAMP stack) --
          \url{https://laragon.org/}
\end{enumerate}

% ──────────────────────────────────────────────────────────────────────────────
\section{Initial Setup (First Time Only)}
% ──────────────────────────────────────────────────────────────────────────────

\subsection{Step 1 -- Clone the Repository into Laragon}

Open your terminal and navigate to your Laragon \texttt{www} directory (e.g., \texttt{C:\textbackslash laragon\textbackslash www}):

\begin{lstlisting}
cd C:\laragon\www
git clone https://github.com/Q-UW9/ro.uvt.ri.git
cd ro.uvt.ri
\end{lstlisting}

\subsection{Step 2 -- Set Up WordPress locally}

Because the WordPress core is ignored in Git (to keep the repo clean), you need to set up a fresh WordPress install inside the project:

\begin{enumerate}[leftmargin=1.5em]
    \item Download WordPress from \url{https://wordpress.org/download/}.
    \item Extract the contents. Move the files into \texttt{ro.uvt.ri\textbackslash wordpress\textbackslash ro.uvt.ri}.
    \item Open Laragon and start all services (\textbf{Start All}).
    \item Open Laragon's database manager (e.g., HeidiSQL) and create a database named \texttt{ro\_uvt\_ri}.
    \item Go to \texttt{wordpress/ro.uvt.ri/} and rename \texttt{wp-config-sample.php} to \texttt{wp-config.php}.
    \item Update the database credentials in \texttt{wp-config.php}:
\end{enumerate}

\begin{lstlisting}
define( 'DB_NAME',     'ro_uvt_ri' );
define( 'DB_USER',     'root' );
define( 'DB_PASSWORD', '' ); // Default Laragon password is empty
define( 'DB_HOST',     'localhost' );
\end{lstlisting}

\begin{enumerate}[leftmargin=1.5em, resume]
    \item In your browser, go to \texttt{http://ro.uvt.ri.test/wordpress/ro.uvt.ri/} and complete the WordPress installation.
\end{enumerate}

\subsection{Step 3 -- Install Frontend Dependencies}

Open a terminal in the \texttt{frontend} folder:

\begin{lstlisting}
cd frontend
npm install
\end{lstlisting}

\subsection{Step 4 -- Configure the API Connection}

Create a \texttt{.env.local} file inside the \texttt{frontend/} directory:

\begin{lstlisting}
# frontend/.env.local
VITE_WP_API_URL=http://ro.uvt.ri.test/wordpress/ro.uvt.ri/wp-json
\end{lstlisting}

\begin{warningbox}
The \texttt{.env.local} file is \textbf{not tracked by Git} (it is listed
in \texttt{.gitignore}). Each developer must create their own copy.
\end{warningbox}

\subsection{Step 5 -- Start the React Development Server}

\begin{lstlisting}
npm run dev
\end{lstlisting}

The React application will start at \texttt{http://localhost:5173}. \\
You can now open both URLs side by side:

\begin{itemize}[leftmargin=1.5em]
    \item \textbf{Backend (CMS):} \texttt{http://ro.uvt.ri.test/wordpress/ro.uvt.ri/wp-admin/}
    \item \textbf{Frontend Website:} \texttt{http://localhost:5173}
\end{itemize}

% ──────────────────────────────────────────────────────────────────────────────
\section{Alternative: WordPress via wp-env (Docker)}
\label{sec:alt-wp}
% ──────────────────────────────────────────────────────────────────────────────

If you have Docker Desktop installed, you can use \texttt{wp-env} to manage WordPress instead of Laragon.

\begin{lstlisting}
npm -g install @wordpress/env
wp-env start
\end{lstlisting}

This will spin up a local WordPress instance on \texttt{http://localhost:8888}. Update your frontend \texttt{.env.local} to use \texttt{http://localhost:8888/wp-json}.

% ──────────────────────────────────────────────────────────────────────────────
\section{Daily Workflow}
% ──────────────────────────────────────────────────────────────────────────────

\subsection{Starting a Work Session}

\begin{lstlisting}
# 1. Pull the latest changes
git pull origin main

# 2. Make sure Laragon is running
# (Click "Start All" in the Laragon control panel)

# 3. Start the React frontend
cd frontend
npm run dev
\end{lstlisting}

\subsection{Creating a Feature Branch}

Always work on a dedicated branch -- never commit directly to \texttt{main}:

\begin{lstlisting}
# Create and switch to a new branch
git checkout -b feature/your-feature-name

# ... make your changes ...

# Stage and commit
git add .
git commit -m "feat: describe your change"

# Push the branch to GitHub
git push origin feature/your-feature-name
\end{lstlisting}

Then open a \textbf{Pull Request} on GitHub to merge into \texttt{main}.

\subsection{Keeping Up to Date}

Before starting new work, always synchronize with the latest \texttt{main}:

\begin{lstlisting}
# Switch to main and pull the latest
git checkout main
git pull origin main

# Switch back to your feature branch and rebase
git checkout feature/your-feature-name
git rebase main
\end{lstlisting}

\begin{tipbox}
If new dependencies were added by a teammate, you may need to re-run
\texttt{npm install} inside the \texttt{frontend/} directory after
pulling.
\end{tipbox}

% ──────────────────────────────────────────────────────────────────────────────
\section{Useful Commands Reference}
% ──────────────────────────────────────────────────────────────────────────────

\begin{center}
\begin{tabular}{|l|l|}
\hline
\textbf{Command} & \textbf{Description} \\
\hline\hline
\texttt{cd frontend \&\& npm install} & Install frontend dependencies \\
\texttt{cd frontend \&\& npm run dev} & Start React dev server (port 5173) \\
\texttt{cd frontend \&\& npm run build} & Build production bundle \\
\texttt{cd frontend \&\& npm run lint} & Run ESLint checks \\
\hline
\texttt{git pull origin main} & Sync with latest remote changes \\
\texttt{git rebase main} & Rebase current branch onto main \\
\texttt{git stash} & Temporarily shelve uncommitted work \\
\hline
\end{tabular}
\end{center}

% ──────────────────────────────────────────────────────────────────────────────
\section{Troubleshooting}
% ──────────────────────────────────────────────────────────────────────────────

\subsection{``Permission Denied'' on Windows Git Add}

If Git reports \texttt{dubious ownership} errors or fails to index files:
\begin{lstlisting}
git config --global --add safe.directory 'C:/laragon/www/ro.uvt.ri'
\end{lstlisting}
Also ensure Laragon is stopped when performing heavy Git operations.

\subsection{CORS Errors When Fetching from WordPress}

Install and activate the
\href{https://wordpress.org/plugins/wp-graphql/}{WP GraphQL} or a CORS
plugin on your local WordPress instance. Alternatively, configure the
Vite proxy in \texttt{frontend/vite.config.ts} to point to \texttt{http://ro.uvt.ri.test}.

% ──────────────────────────────────────────────────────────────────────────────
\section{Summary}
% ──────────────────────────────────────────────────────────────────────────────

\begin{center}
\fbox{\parbox{0.85\textwidth}{
\centering
\textbf{Clone to Laragon} $\rightarrow$
\textbf{Setup WP locally} $\rightarrow$
\textbf{npm install} $\rightarrow$
\textbf{npm run dev} $\rightarrow$
\textbf{Start coding!}
}}
\end{center}

\vspace{1em}

For questions or issues, open an issue on the GitHub repository: \\
\url{https://github.com/Q-UW9/ro.uvt.ri}

\end{document}
